\subsection{绝对值的情形}

定理 2 （绝对值的逼近定理）. 设 \(f_{i}\) ( \(1 \leqslant i \leqslant n\) ) 是同一域 \(\mathbf{K}\) 上不是非真的且两两不等价的绝对值。设 \(a_{i}\) ( \(1 \leqslant i \leqslant n\) ) 是 \(\mathbf{K}\) 的元素，\(\varepsilon\) 是实数 \(>0\) 。那么存在 \(x \in \mathbf{K}\) 使得 \(f_{i}(x - a_{i}) \leqslant \varepsilon\) 对一切 \(i\) 成立。

用 \(K_{i}\) 表示赋以由 \(f_{i}\) 定义的拓扑的域 K 。要证的结果等价于：在乘积 \(P = K_{1} \times \cdots \times K_{n}\) 中，对角线 D 的闭包 \(\overline{D}\) 等于 P。对 n = 1 这是显然的。设这一点对 k 个绝对值（k < n）的情形已经确立。

首先证明，对 \(2 \leqslant h \leqslant n\) ，存在 K 的元素 \(x_{h}\) 使得 \(f_{1}(x_{h}) < 1\) 、\(f_{2}(x_{h}) > 1\) 且 \(f_{i}(x_{h}) \neq 1\) 对 \(3 \leqslant i \leqslant h\) 成立。我们对 h 作归纳论证。若 h = 2 ，这由 \(f_{1}\) 与 \(f_{2}\) 不等价这一事实得出。于是设 \(x_{h-1}\) 的存在性已证，而来证明 \(x_{h}\) 的存在性。若 \(f_{h}(x_{h-1}) \neq 1\) ，可取 \(x_{h} = x_{h-1}\) ；若 \(f_{h}(x_{h-1}) = 1\) ，选取 \(z \in K^{*}\) 使得 \(f_{h}(z) \neq 1\) ，则对充分大的 s ，\(x_{h} = z(x_{h-1})^{s}\) 解决问题。这样就证明了 \(x_{h}\) 的存在性。

当整数 \(q\) 趋于无穷时，\(f_{1}(x_{n}^{q})\) 趋于 0，\(f_{2}(x_{n}^{q})\) 趋于 \(+\infty\) ，且 \(f_{i}(x_{n}^{q})\) 趋于 0 或 \(+\infty\) 对 \(i \geqslant 3\) 成立。记 \(y_{q} = x_{n}^{q}(1 + x_{n}^{q})^{-1}\) ，

\[
1 - y _ {q} = (1 + x _ {n} ^ {q}) ^ {- 1};
\]

于是序列 \((y_q)\) 在 \(\mathbf{K}_1\) 中趋于 0，在 \(\mathbf{K}_2\) 中趋于 1，且在 \(\mathbf{K}_i\) 中趋于 0 或 1（对 \(i \geqslant 3\) ）。改变诸 \(\mathbf{K}_i\) 的编号，因此可以设存在整数 \(r\) ( \(1 \leqslant r < n\) ) 使得 \(\overline{\mathbf{D}}\) 包含点 \((e_1, \ldots, e_n)\) ，其中 \(e_i = 1\) 对 \(1 \leqslant i \leqslant r\) ，\(e_i = 0\) 对 \(r + 1 \leqslant i \leqslant n\) 。现在，\(\overline{\mathbf{D}}\) 是 P 的向量子 K-空间。于是 \(\overline{\mathbf{D}}\) 包含对角线 \(\mathbf{D}'\) 与 \(\mathbf{D}''\) ，它们分别是

\[
\mathbf {P} ^ {\prime} = \mathbf {K} _ {1} \times \dots \times \mathbf {K} _ {r}
\]

与 \(P'' = K_{r+1} \times \cdots \times K_n\) 的对角线。由归纳假设，\(P' = \overline{D}'\) 且 \(P'' = \overline{D}''\) 。于是 \(\overline{D} = P\) 。

